% from main.tex
%%%%%%%%%%  scan idx 153  |  folio 140  %%%%%%%%%%
\section*{6. Rang d'un complexe parfait}

\textbf{6.0.} Soit \((X,\mathcal O_X)\) un espace annelé en anneaux locaux. Il est bien connu que l'on peut attacher à tout faisceau localement libre de type fini \(L\) sur un ouvert \(U\) de \(X\) une fonction \(\rg_U(L)\) localement constante sur \(U\) et à valeurs entières, qu'on appelle le \emph{rang} de \(L\), la famille des fonctions \(\rg_U(L)\) étant caractérisée par les propriétés suivantes :
\begin{enumerate}[label=(\roman*)]
\item si \(V\subset U\) est une inclusion d'ouverts, on a \(\rg_V(L|V)=\rg_U(L)|V\);
\item pour toute suite exacte de \(\mathcal O_U\)-Modules localement libres de type fini
\[
  0\longrightarrow L'\longrightarrow L\longrightarrow L''\longrightarrow 0,
\]
on a \(\rg_U(L)=\rg_U(L')+\rg_U(L'')\);
\item \(\rg_X(\mathcal O_X)=1\).
\end{enumerate}
On se propose, dans ce numéro, d'étendre la notion de rang aux complexes parfaits.

\textbf{Définition 6.1} (cf. SGA 5 VIII 2). Soient \(\mathcal C\) une catégorie additive (resp. triangulée) et \(F\) un groupe abélien. On dit qu'une fonction \(f:\Ob(\mathcal C)\longrightarrow F\) est une fonction additive de \(\mathcal C\) dans \(F\) si pour tous \(L,L',L''\in\Ob(\mathcal C)\) tels que l'on ait \(L\simeq L'+L''\) (resp. un triangle distingué
\[
\begin{tikzcd}
& L'' \arrow[dl, dashed] & \\
L' \arrow[rr] && L \arrow[ul]
\end{tikzcd}
\]
) on a \(f(L)=f(L')+f(L'')\).

Les fonctions additives de \(\mathcal C\) dans \(F\) forment un groupe abélien, noté \(\Add(\mathcal C,F)\).

\textbf{6.2.}
